\chapter{Обзор микроконтроллера MIK32}
\label{cha:mik32}

\section{Архитектура и технические характеристики MIK32}

Микроконтроллер MIK32 представляет собой 32-битное RISC-V решение отечественной разработки компании НПК <<Микрон>>, ориентированное на задачи промышленной автоматизации и интернета вещей \cite{mik32_technical_spec, cnxsoft2021mik32}. Данное решение демонстрирует возможности открытой архитектуры RISC-V в создании конкурентоспособных микроконтроллеров, сопоставимых по функциональности с популярными зарубежными аналогами серии STM32L0 \cite{electronicslab2021mik32}.

Архитектура MIK32 сочетает современные подходы к проектированию встраиваемых систем с требованиями энергоэффективности и информационной безопасности. Микроконтроллер позиционируется как решение для импортозамещения в критически важных областях применения \cite{mik32_technical_spec}.

Основные технические характеристики микроконтроллера представлены в таблице \ref{tab:MIK32_main_features}.

\begin{center}
	\begin{longtable}{|p{0.34\textwidth}|p{0.60\textwidth}|}
		\caption{Основные характеристики микроконтроллера MIK32}
		\label{tab:MIK32_main_features}
		\\ \hline
		Характеристика & Значение \\
		\hline \endfirsthead
		\subcaption{Продолжение таблицы~\ref{tab:MIK32_main_features}} \\
		\hline
		Характеристика & Значение \\
		\hline \endhead
		\hline \subcaption{Продолжение на след. стр.} \endfoot
		\hline \endlastfoot
		
		Процессорное ядро & RISC-V 32-бита RV32IMC (SCR1, Syntacore) \\
		\hline
		Максимальная тактовая частота & 32 МГц \\
		\hline
		ОЗУ & 16 КБ \\
		\hline
		Энергонезависимая память & 8 КБ EEPROM, 256 бит OTP \\
		\hline
		Внешняя память & До 16 МБ через QSPI Flash (Single/Dual/Quad SPI) \\
		\hline
		DMA & 4 канала, поддержка различных режимов передачи \\
		\hline
		Таймеры & 3×32-битных, 3×16-битных с поддержкой ШИМ \\
		\hline
		Интерфейсы связи & 2×UART, 2×SPI, 2×I2C \\
		\hline
		GPIO & 40 линий (конфигурируемые) \\
		\hline
		АЦП & 12 бит, 8 каналов, до 1 МГц \\
		\hline
		ЦАП & 12 бит, 2 канала, до 1 МГц \\
		\hline
		Криптография & CRC32, AES-128, ГОСТ Р 34.12-2015 \\
		\hline
		Питание & 3,3 В ±10\% \\
		\hline
		Температурный диапазон & –45°C ... +85°C \\
		\hline
		Корпус & QFN-64 \\
		\hline
	\end{longtable}
\end{center}

В основе MIK32 лежит ядро SCR1 компании Syntacore, реализующее архитектуру RISC-V с базовым набором инструкций RV32IMC \cite{mik32_technical_spec}. Архитектура RISC-V предоставляет гибкие и масштабируемые механизмы взаимодействия с аппаратным окружением, которые охватывают работу с памятью, регистрами, периферийными устройствами, а также обработку прерываний и исключений. Эти функции критически важны для создания эффективных системных приложений, особенно в операционных системах и встраиваемых устройствах.

Конвейерная организация ядра включает следующие основные компоненты:

\begin{itemize}
	\item \textbf{Модуль выборки инструкций (IFU)} — обеспечивает получение команд из памяти через интерфейс AHB-Lite, оптимизирован для работы с 32-битными и 16-битными инструкциями расширения C;
	\item \textbf{Модуль декодирования (IDU)} — выполняет анализ инструкций и определение необходимых ресурсов для их выполнения;
	\item \textbf{Исполнительный модуль (EXU)} — реализует арифметико-логические операции, включая аппаратную поддержку умножения и деления благодаря расширению M;
	\item \textbf{Модуль обращения к памяти (LSU)} — организует доступ к данным через интерфейс AHB-Lite с поддержкой различных типов операций загрузки и сохранения.
\end{itemize}

Система управления включает программируемый контроллер прерываний (IPIC), поддерживающий до 32 векторных прерываний с конфигурируемыми приоритетами \cite{mik32_technical_spec}. Встроенные средства отладки реализованы через JTAG-интерфейс, совместимый со стандартом IEEE 1149.1-2013.

Архитектура ядра представлена на рисунке \ref{fig:scr1_arch}.

\begin{figure}[ht]
	\centering
	\includegraphics[width=0.7\textwidth]{inc/img/scr1-arch}
	\caption{Архитектура ядра SCR1 микроконтроллера MIK32}
	\label{fig:scr1_arch}
\end{figure}

Таким образом, проведенный обзор архитектуры и технических характеристик микроконтроллера MIK32 позволяет сделать следующие выводы:
\begin{itemize}
	\item Архитектурная состоятельность — MIK32 демонстрирует успешную адаптацию открытой архитектуры RISC-V для задач встраиваемых систем с обеспечением конкурентоспособных характеристик.
	\item Функциональная полнота — набор периферийных устройств и интерфейсов микроконтроллера соответствует требованиям современных приложений промышленной автоматизации и интернета вещей.
	\item Стратегическая значимость — реализация отечественного микроконтроллера с поддержкой криптографических стандартов ГОСТ создает основу для решения задач импортозамещения в критически важных областях.
\end{itemize}

Полученные результаты создают необходимую базу для перехода к рассмотрению программного обеспечения и инструментов разработки для микроконтроллера MIK32 в следующей главе.

% Определения и сокращения
\Abbrev{RISC-V}{Reduced Instruction Set Computer V ""--- открытая архитектура процессоров с сокращённым набором команд}
\Define{SCR1}{32-битное RISC-V ядро компании Syntacore с поддержкой расширений RV32IMC}
\Abbrev{SPI}{Serial Peripheral Interface — последовательный периферийный интерфейс}
\Define{АЦП}{аналого-цифровой преобразователь для преобразования непрерывных сигналов в цифровой код}
\Define{ЦАП}{цифро-аналоговый преобразователь для преобразования цифровых данных в аналоговый сигнал}
\Abbrev{DMA}{Direct Memory Access — прямой доступ к памяти без участия процессора}
\Abbrev{EPIC}{Enhanced Programmable Interrupt Controller — расширенный программируемый контроллер прерываний}
